\chapter{Reproduction}
\label{app:reproduction}

\paragraph{Data and code availability.}  The raw sources are public
(\cref{sec:data:collection}): the order-book archive on Kaggle, the Polygon blockchain and
the Polymarket metadata API.  Rebuilding every derived dataset from them is possible but
long, so the datasets that the models read are deposited as they are: the market table
(175{,}292 rows), the hourly transaction partitions with the market and the token
orientation of every fill (11~GB), the order-book tensors (27~GB), the hand-made features
(7.7~GB), the cross-fitted tree log-probabilities (6.6~GB), the foundation-model forecasts
(0.7~GB), the news and sibling features and the identifier mappings, about 55~GB in all,
together with the trained checkpoints and the reports from which every table is taken
\todo{URL of the data repository once chosen}.  With them, every model of the thesis can be
retrained and evaluated on the same inputs without touching the raw streams.  The code is
a Python package (\path{polymarket_gnn}) managed with \path{uv}, with the data pipeline in
DuckDB and NumPy and the models in PyTorch \todo{URL of the code repository}; the raw data
under \path{data/raw/} are never modified and every derived dataset lives under
\path{data/derived/}.

\paragraph{Steps.}  The steps that produce the numbers of the thesis are, in order:

\begin{enumerate}[nosep]
  \item \textbf{Collection.}  \path{scripts/preprocessing/00_legacy_raw_downloads/}: the
  Gamma API pager for the markets, the markets-by-token completion, and the Polygon block
  reader for the fills, a modified \path{poly-trade-scan} \cite{martkir2026polytradescan}
  (\cref{sec:data:collection}).
  \item \textbf{Market tables.}  \path{sql/markets/00}--\path{09} run by
  \path{scripts/preprocessing/00_build_markets_tables.py}: the Gamma rows flattened into
  one table per market, every fill joined to its market through the two outcome tokens,
  the union of the market identifiers seen in the five sources and the completion of the
  missing ones by condition identifier, with the three API scrapers that run between the
  steps and the expected row counts in \path{sql/markets/README.md}.
  \item \textbf{Preprocessing.}  \path{scripts/preprocessing/01}--\path{06}: multi-day
  aggregates, hourly order-book partitions with the 52-feature vectors, deltas, text
  embeddings, node features, wallet features; \path{precompute_tensors.py} for the
  per-block tensors read by the stateful model.
  \item \textbf{Causal baselines.}  \path{causal_baselines.py} \path{--horizon} \path{H} writes
  \path{baseline_results_<H>.json} with every rule on all and deduplicated samples.
  \item \textbf{Stateful runs.}  \path{run_hand_runs_v4.sh} \path{<device>} \path{<H>} with the
  environment variables of \cref{app:runs} (\path{EPOCHS}, \path{TREE_DIR},
  \path{TREE_OFFSET}, \path{ENC}, \path{W}, \path{SEED}); each run writes its
  validation predictions.
  \item \textbf{Hand features and the control.}  \path{hand_features.py} builds the 61
  features for every tensor row; \path{scripts/data_audit/feature_audit.py} \path{H} \path{<label>
} \path{} \path{<label>} \path{10} \path{<out-dir>} builds the control on the samples of the run \path{<label>}, group
  by group, with the univariate diagnostics and the importances;
  \path{control_proba.py} \path{--horizon} refits G5 for the class probabilities.
  \item \textbf{Data-mining study.}  \path{scripts/data_audit/datamining.py} \path{--horizon}
  \path{H} reads the feature parquets of the previous step and writes, per horizon, the PCA,
  correlation, random-projection, MDS, UMAP and t-SNE outputs, the LDA, logistic, kNN, tree
  and mutual-information results and the market clustering of \cref{ch:structure}, with a
  \path{report.md}; CPU only, about twenty-five minutes per horizon.
  \item \textbf{Transaction graph.}  \path{scripts/data_audit/graph_animation.py}
  \path{--stage} \path{E,S} builds the edge events and the statistics of \cref{sec:expl:graph}
  (degree distributions, components, $k$-cores, hourly and daily series);
  \path{graph_layout_anchor.py} (with \path{--cat-weight 0.5} for the second map) computes
  the SVD and UMAP layout, \path{graph_islands.py} the HDBSCAN islands, \path{--stage F,V} the
  frames and the video, and \path{graph_summary_plots.py} the summary figure; CPU only.
  \item \textbf{Cross-fitted trees and the residual model.}  \path{tree_features.py}
  writes the day-cross-fitted log-probabilities; the residual run is step 4 with
  \path{TREE_DIR} and \path{TREE_OFFSET=1}.
  \item \textbf{Tabular networks.}  \path{scripts/data_audit/tabular_net.py} \path{--horizon
} \path{} \path{--variants} \path{tabm} \path{--seeds} \path{0,1,2}.
  \item \textbf{Comparisons.}  \path{hand_vs_lgb_report.py} \path{<label>} \path{--horizon} \path{--prior
} \path{} \path{--crossfit} (raw, prior-corrected and one-sided offsets, by subset and category);
  \path{scripts/data_audit/offset_fair.py} \path{--horizon} \path{--deep} \path{<labels>} \path{--tab} \path{<probas>
} \path{} \path{--grid-max} \path{2.0} \path{--grid-step} \path{0.2} (the same calibration for every model);
  \path{analyze_predictions.py} (per-market breakdown);
  \path{scripts/data_audit/leadlag.py} (sibling lead--lag);
  \path{news_features.py}, \path{sibling_features.py} and
  \path{scripts/data_audit/news_audit.py} (external features in the control).
  \item \textbf{Reports.}  Every table of the thesis is in
  \path{reports/experiments/<study>/}, with the run labels in the file names; the
  development log (\path{DEVELOPMENT_LOG.md}) records every run with its date, device
  and duration.
\end{enumerate}

Unit tests cover the label construction, the mid-price index, the hand-feature index and
the news index (\path{tests/}).  The thesis is built with \path{latexmk} from the
\path{thesis/} directory; the bibliography is tagged by the kind of verification of each
entry (\path{[L]} a PDF in \path{literature/}, \path{[V]} checked online, \path{[A]}
abstract only).
